\section{Illustrative examples}
\label{sec:results}

This section demonstrates \emph{how the tool is used} and the \emph{shape of its output}
through four short example analyses. The numbers are real \tool runs, included to show that
the instrument works and what it produces---not as standalone scientific findings; the
substantive studies that \tool supports are reported in the authors' companion work on
graph-based spear-phishing detection (in preparation). Each example reproduces from a bundled
CSV via the matching \texttt{experiments} command.

\subsection{Running the detector panel}
A single command builds a labelled scenario and scores the whole panel:

\begin{lstlisting}
$ python -m cascadebench demo --graph enron --fanout 8
graph: enron, 150 nodes
scenario: 4213 events, 37 attacks, 12% infected
  1-hop            AUC=0.65  R@1%=0.02
  hand-context     AUC=0.77  R@1%=0.04
  temporal-GNN     AUC=0.85  R@1%=0.07
  ...                                  # abridged: the full run prints all 8 detectors
\end{lstlisting}

Running the panel under two attacker regimes (Table~\ref{tab:panel}) illustrates the kind of
comparison \tool is built to surface: because the attacker is programmable, the same panel can
be scored against a naive and a stealthy adversary, and the tool reports both side by side with
confidence intervals. Whether and how the ranking changes between regimes is exactly the
question a user studies with the tool---a comparison a fixed corpus cannot produce.

\begin{table}[H]
\centering
\caption{Example \tool output: detector panel AUC under two attacker regimes (synthetic
organization, $\ge$20 seeds, 95\% CI half-widths $\le 0.02$). Shown to illustrate the tool's
output format, not as a benchmark result.}
\label{tab:panel}
\small
\begin{tabular}{@{}lcc@{}}
\toprule
Detector & AUC (naive) & AUC (stealthy) \\
\midrule
1-hop baseline                       & 0.70 & 0.63 \\
COMPA \cite{egele2013compa}          & 0.83 & 0.67 \\
hopper \cite{ho2021hopper}           & 0.81 & 0.71 \\
anomaly forest                       & 0.63 & 0.72 \\
static GCN \cite{kipf2017semisupervised} & 0.80 & 0.84 \\
hand-crafted context                 & 0.75 & 0.69 \\
temporal GNN \cite{rossi2020temporal} & 0.87 & 0.75 \\
temporal GNN + attention \cite{xu2020inductive} & 0.87 & 0.75 \\
\bottomrule
\end{tabular}
\end{table}

\subsection{Auditing for leakage}
The \texttt{leak-audit} command applies the time-shuffle control and reports, as a first-class
number, how much of a detector's advantage survives permuting the timestamps
(Table~\ref{tab:leak}). This is the diagnostic \tool contributes: a user can check whether a
``temporal'' result rests on genuine dynamics or on a timing confounder, and report the gap
rather than a single optimistic score.

\begin{table}[H]
\centering
\caption{Example \tool output: the time-shuffle leak audit (Enron graph, 8 seeds). The tool
reports the shuffled and unshuffled scores so the confound-driven portion is explicit.}
\label{tab:leak}
\small
\begin{tabular}{@{}lcc@{}}
\toprule
Setting & AUC & Recall@FPR=1\% \\
\midrule
temporal GNN                 & 0.85 & 0.072 \\
temporal GNN (time-shuffled) & 0.77 & 0.049 \\
1-hop baseline               & 0.65 & 0.022 \\
\bottomrule
\end{tabular}
\end{table}

\subsection{Sweeping the attacker}
\texttt{pareto} and \texttt{adaptive} sweep the attacker's reach and record a curve rather than
a point. In the example run, as \texttt{fanout} falls from 8 to 2 (infected fraction
$0.92\!\to\!0.40$) the per-detector AUC curves cross, so the tool returns a reach--evasion
\emph{frontier} for each detector. This is the output a user needs to reason about robustness
against an adaptive opponent, and it is what the companion study consumes; the tool's job is to
produce it reproducibly.

\subsection{Transfer analysis}
The \texttt{transferability} command optimizes an attack against \emph{each} detector and
evaluates every optimized attack against \emph{all} detectors, emitting an
attack$\times$detector matrix (\texttt{cb\_transfer.csv}); off-diagonal entries quantify how a
tuned attack transfers. Because the same interface loads both the synthetic organization and
the real Enron graph, the identical protocol runs on synthetic and real structure, so a user
can inspect ranking concordance across graph sources. \tool thus also serves as a check on its
own external validity---a capability a single-corpus study cannot offer.